% Srikar's section. Insert immediately before \end{document}.
% Numerical results remain blank pending consistent upstream regeneration.
\clearpage
\section{Circuit-Level Validation: Parity Checker}
\label{sec:circuit}

\subsection{Aim}
This part extends the gate-level method to a complete circuit. The aim is
to estimate its static supply leakage by combining predictions for its
individual gates, and compare the estimate with a transistor-level SPICE
simulation of the complete circuit.

\subsection{Circuit construction}
The circuit checks eight data bits, $d_0,\ldots,d_7$, together with an
even-parity bit $p$. Its output is
\begin{equation}
e=d_0\oplus d_1\oplus\cdots\oplus d_7\oplus p.
\end{equation}
An output of 0 indicates valid even parity; an output of 1 indicates a
parity error. There are nine primary inputs and $2^9=512$ input vectors.
The vector order used throughout is $d_0,d_1,\ldots,d_7,p$.

The eight data bits are combined using a balanced XOR tree. Its output is
then XORed with the parity bit. Each of the eight XOR operations is built
from four NAND2 gates:
\begin{align}
t_1&=\operatorname{NAND}(a,b), &
t_2&=\operatorname{NAND}(a,t_1),\\
t_3&=\operatorname{NAND}(b,t_1), &
a\oplus b&=\operatorname{NAND}(t_2,t_3).
\end{align}
The complete circuit therefore contains 32 NAND2 gates and 128 transistors.
Each NAND2 uses the same topology and minimum transistor sizing as Part 3:
$W=L=22$\,nm. The intended consistent device setup uses drain and source
areas $AD=AS=W(2L)$. The transistor netlist contains no ideal logic gates
and no added output load.

One ordered list of gate connections is used for both Boolean evaluation
and netlist generation. NAND inputs retain their order: input $a$ drives
the upper NMOS in the stack and input $b$ drives the lower NMOS. This matters
because states \texttt{01} and \texttt{10} need not have the same leakage.

\subsection{PVT conditions and SPICE reference}
The evaluation uses the same 200 PVT corners as Part 3. Each corner specifies
the supply voltage, temperature, and separate NMOS and PMOS process
parameters. Evaluating every input vector at every corner gives
$512\times200=102{,}400$ circuit cases. These are 200 distinct PVT draws,
each reused for all 512 vectors.

The reference leakage is the current drawn from the complete circuit's
shared supply, $I_{\mathrm{SPICE}}=-i(\texttt{vsupply})$. Current supplied by
the separate primary-input sources is excluded from this main definition.
The simulator records the supply current, primary-input currents, and all
32 NAND output voltages. A DC sweep enumerates the binary inputs using
sources that hold each input at exactly 0 or $V_{DD}$. The circuit itself
remains a transistor-level circuit. Independent operating-point simulations
provide a cross-check of this sweep method.

\subsection{Hierarchical estimate from the transistor models}
The circuit estimate reuses the Part 3 NAND2 predictions. At each corner,
a lookup table contains the four ordered NAND input states, giving
$4\times200=800$ entries. Each entry contains the gate's predicted supply
leakage and the signed gate currents at its two input pins. The supply
leakage uses the KCL method described in Section~\ref{sec:combine}.
For a NAND input pin, its current is the sum of the gate-terminal currents
of the PMOS and NMOS connected to that pin.

For each nine-bit input vector, Boolean evaluation determines the input
state of every NAND. The corresponding lookup entry is selected for each
gate and the contributions are combined. The circuit estimate does not
use node voltages from the complete-circuit SPICE simulation. However,
the isolated-gate lookup inherits the SPICE-recorded transistor voltages
used by Part 3; it is therefore a hierarchical approximation based on
characterised gates.

\subsection{Combining the gate contributions}
\paragraph{Plain gate sum (baseline).}
The prescribed hierarchical estimate adds the supply leakages of the
32 isolated NAND gates:
\begin{equation}
\widehat I_{\mathrm{plain}}=\sum_{j=1}^{32}\widehat I_{\mathrm{supply},j}.
\end{equation}
This baseline is retained in the comparison.

\paragraph{Correction for internally driven high inputs.}
An isolated gate's input is powered by an external voltage source. Inside
the circuit, an internal high input is instead driven by another gate,
which must supply its input gate current. The plain sum omits this
contribution. Under the ideal-rail approximation, the corrected estimate is
\begin{equation}
\widehat I_{\mathrm{corrected}}=
\widehat I_{\mathrm{plain}}+
\sum_{(j,q)\in\mathcal H_{\mathrm{internal}}}\widehat I_{g,jq},
\end{equation}
where $\mathcal H_{\mathrm{internal}}$ contains the receiving input pins
that are high and driven by another gate. Currents are signed, positive
into the receiving pin; their absolute values are not used. Each receiving
pin is counted once, including when a driver has multiple loads.

Primary-input currents remain excluded because their sources are separate
from the circuit supply. A separate variant adds the currents at high
primary inputs to both the prediction and the SPICE reference, so that the
two sides still use the same leakage definition. The correction changes
current accounting; it does not solve the circuit's coupled analog voltages.

\subsection{Validation procedure}
The implementation includes the following checks:
\begin{itemize}
\item Exhaustive Boolean evaluation of all 512 vectors, including checking
      that flipping the parity bit flips the error output.
\item Structural checks for an ordered, acyclic graph with 32 NAND2 gates
      and a generated netlist containing 128 transistors.
\item Checks of every simulated NAND output against its Boolean value,
      primary-input voltages against their required rails, and complete
      coverage of the corner--vector pairs.
\item Tests of signed input-current ownership, agreement between batched
      and individual aggregation, and rejection of missing or incorrect
      simulation measurements.
\item An independent operating-point versus DC-sweep comparison.
\end{itemize}
\noindent\textit{Final validation outcomes and simulator version: pending
the run with consistently regenerated datasets and models.}

\subsection{Results}
For each circuit case, the relative error and signed relative error are
\begin{equation}
r=\frac{|\widehat I-I_{\mathrm{SPICE}}|}{|I_{\mathrm{SPICE}}|},
\qquad
b=\frac{\widehat I-I_{\mathrm{SPICE}}}{I_{\mathrm{SPICE}}}.
\end{equation}
The comparison reports median, 95th-percentile and mean relative error;
bias is the median signed relative error. Under the project's convention,
mean accuracy is $100(1-\operatorname{mean}(r))\%$.

\begin{table}[H]
\centering\small
\begin{tabular}{lrrrr}
\toprule
Estimate & Bias (\%) & Median error (\%) & p95 (\%) & Mean (\%)\\
\midrule
Plain gate sum & & & & \\
Corrected supply estimate & & & & \\
Including primary high inputs & & & & \\
\bottomrule
\end{tabular}
\caption{Circuit comparison. Numerical cells are intentionally blank pending
the final consistent run.}
\label{tab:circuitresults}
\end{table}

\begin{table}[H]
\centering
\begin{tabular}{lr}
\toprule
Quantity (supply-only comparison) & Value\\
\midrule
Median SPICE leakage ($\mu$A) & \\
Median corrected ML leakage ($\mu$A) & \\
Maximum corrected relative error (\%) & \\
Corrected mean accuracy (\%) & \\
\bottomrule
\end{tabular}
\caption{Additional circuit metrics, to be filled after the final run.}
\end{table}

\begin{figure}[H]
\centering
\fbox{\parbox[c][35mm][c]{0.85\linewidth}{\centering
Pending final circuit parity plot: plain and corrected ML estimates
against complete-circuit SPICE, with logarithmic axes.}}
\caption{Placeholder for the final circuit leakage comparison.}
\label{fig:circuitparity}
\end{figure}

\begin{figure}[H]
\centering
\fbox{\parbox[c][35mm][c]{0.85\linewidth}{\centering
Pending final plot of corrected relative error against temperature.}}
\caption{Placeholder for the final temperature-dependent error plot.}
\label{fig:circuittemperature}
\end{figure}

\subsection{Worked calculation}
For vector \texttt{000000000} at corner 0, Boolean evaluation gives eight
NAND gates in state \texttt{00}, sixteen in \texttt{01}, none in
\texttt{10}, and eight in \texttt{11}. Thus,
\begin{equation}
\widehat I_{\mathrm{plain}}=
8\widehat I_{00}+16\widehat I_{01}+8\widehat I_{11}.
\end{equation}
The internal high-input currents are then added and the result is compared
with the complete-circuit supply current.
\begin{table}[H]
\centering
\begin{tabular}{lr}
\toprule
Calculation step & Value\\
\midrule
Plain gate sum ($\mu$A) & \\
Internal high-input correction ($\mu$A) & \\
Corrected estimate ($\mu$A) & \\
Complete-circuit SPICE reference ($\mu$A) & \\
Relative error (\%) & \\
\bottomrule
\end{tabular}
\caption{Worked example; numerical values await the final run.}
\end{table}

\subsection{Error analysis and limitations}
\noindent\textit{Results-dependent discussion is pending: comparison with
gate-level accuracy, worst corners and vectors, temperature dependence,
and whether the mean-accuracy objective is reached.}

The lookup assumes ideal input rails and isolated-gate operating points.
In the complete circuit, fanout and departures from ideal high and low
voltages can change leakage. The input-current correction accounts for
internally supplied input currents but does not remove this approximation.
The transistor models also retain prediction errors.

To distinguish these effects, the same aggregation is also evaluated using
isolated-gate SPICE currents in place of ML currents. Its disagreement with
complete-circuit SPICE measures the error of composing isolated gates;
comparison with the ML estimate helps assess the additional model error.
Absolute errors from these comparisons need not add, because signed errors
can cancel.

Possible improvements include gate models conditioned on analog input
voltages and loading, a coupled node-voltage solution, and improved
transistor training coverage. Circuit evaluation cases must remain outside
the training data.

\subsection{Reproducibility}
The circuit construction, simulation, comparison and validation scripts
are provided in \texttt{part4\_circuit/}. Once the consistent upstream
artifacts are available, \texttt{python run\_all.py --part4-only} runs the
Part 4 tests, circuit simulation workflow and comparison. The separate
\texttt{part4\_circuit/cross\_check.py} script performs the independent
operating-point/DC check. Final tables, plots and validation records will
be taken from that consistent run.
